
\section{Introduction}
Given a functional $\calF:\calP(\R^d)\to\R$, its gradient-flow dynamics depend on the choice of geometry on the underlying space of measures. One of the most commonly used geometries is the Wasserstein-2 geometry, induced by the Wasserstein-2 distance. Its associated Onsager operator, which is the inverse of the Riemannian metric tensor, is given by
$\mathbb{K}_W(\rho): T_\rho^* \mathcal{M}^{+} \rightarrow T_\rho \mathcal{M}^{+}, \xi \mapsto-\operatorname{div}(\rho \nabla \xi)$.
Hence, the Wasserstein gradient flow of $\calF$ is
\begin{align}
\partial_t \mu_t = \nabla\cdot \left(\mu_t \nabla \calF'[\mu_t]\right). \tag{W-GF}
\end{align}
Similarly, the Fisher--Rao geometry is induced by the Hellinger distance, with associated Onsager operator
$\mathbb{K}_{\mathrm{FR}}(\rho): T_\rho^*\mathcal M^+\rightarrow T_\rho\mathcal M^+, \xi\mapsto \rho , \xi$.
Consequently, the Fisher--Rao gradient flow of $\calF$ is
\begin{align}
\partial_t\mu_t = -\mu_t \left( \calF'[\mu_t] - \int \calF'[\mu_t]  \dd \mu_t \right). \tag{FR-GF}
\end{align}
An interpolation between these two dynamics is given by the well-known Wasserstein--Fisher--Rao gradient flow:
\begin{align}
\partial_t\mu_t = \alpha \nabla\cdot \left(\mu_t \nabla \calF'[\mu_t]\right) - \beta \mu_t \left( \calF'[\mu_t] - \int \calF'[\mu_t]  \dd \mu_t \right). \tag{WFR-GF}
\end{align} 
Recently, a different geometry is considered in \cite{zhu2024inclusive,zhu2024kernel,gladin2024interaction}, which uses the MMD as the geometry. The associated Onsager operator is $\mathbb{K}_{\mathrm{MMD}}(\rho) = \mathbb{K}_{\mathrm{FR}}(\rho) \circ \mathcal{K}_\rho^{-1}$. The Interaction-Force Transport (IFT) gradient flow is described as the following:
\begin{align}
\partial_t\mu_t = -\mathbb{K}_{\mathrm{MMD}}(\mu_t) \calF'[\mu_t] = -\mu_t \mathcal{K}_{\mu_t}^{-1} \calF'[\mu_t] = - \mathcal{K}^{-1} \calF'[\mu_t]. \tag{IFT-GF}
\end{align}
Similarly as WFR-GF, an interpolation of WGF and IFT-GF can also be constructed
\begin{align*}
    \partial_t\mu_t = \alpha \nabla\cdot \left(\mu_t \nabla \calF'[\mu_t]\right) - \beta \mathcal{K}^{-1} \calF'[\mu_t] \tag{WIFT-GF}.
\end{align*}


\section{$\calF=\mmd$ and geometry=$\mmd$}
We focus on MMD due to its interpretation as the worse-case integration error in the unit ball of a RKHS. 
If $\calF$ is the MMD distance $\calF(\cdot)=\mmd^2(\cdot, \pi)$, then $\calF'[\mu](\cdot)=\int k(x, \cdot) \dd(\mu-\pi) = \mathcal{K}(\mu-\pi)$. 
MSIP uses the WFG-GF, but here we focus on the WIFT-GF. 
The main reason is that, \cite{gladin2024interaction} has proved that $\operatorname{MMD}^2\left(\mu_t, \pi\right) \leq e^{-2 \beta t} \cdot \operatorname{MMD}^2 \left(\mu_0, \pi\right)$. 

Following \cite{gladin2024interaction}, a way to simulate the WIFT-GF.  is via an weighted interacting particle system.  
Let $\mu_N^{(t)}=\sum_{i=1}^N w_i^{(t)} \delta_{x_i^{(t)}}$ be a weighted empirical measure at iteration $t$, where the weights do not necessarily sum up to $1$. 
Denote $m_\pi=\int k(x,\cdot)\dd \pi$. 
First evolve the Wasserstein gradient term, 
\begin{align*}
    x_i^{(t+1)}=x_i^{(t)}-\tau \alpha\left[\sum_{j=1}^N w_j^{(t)} \nabla_2 k\left(x_j^{(t)}, x_i^{(t)}\right)-\nabla m_\pi\left(x_i^{(t)}\right)\right].
\end{align*}
Then, evolve the IFT gradient term, by solving the following optimal weights: let $\eta=\beta\tau \alpha^{-1}$, 
\begin{align*}
    w^{(t+1)} = \arg \min _{w \in \R^N}\left\{\frac{1}{2} \mmd^2\left(\sum_{i=1}^N w_i \delta_{x_i^{(t+1)}}, \pi\right)+\frac{1}{2 \eta} \mmd^2\left(\sum_{i=1}^N w_i \delta_{x_i^{(t+1)}}, \mu_N^{(t)}\right)\right\}.
\end{align*}
Interestingly, let $K^{(t+1)}=k(x_{1:N}^{(t+1)},x_{1:N}^{(t+1)})\in\R^{N\times N}$, the weight problem is
\begin{align*}
    \min_w \left\{\frac{1}{2} w^{\top} K^{(t+1)} w - m_\pi(x_{1:N}^{(t+1)})^{\top} w + \frac{1}{2 \eta} \left(w-w^{(t)} \right)^{\top} K^{(t+1)} \left(w-w^{(t)}\right)\right\} ,
\end{align*}
By definition of kernel quadrature, there is $m_\pi(x_{1:N}^{(t+1)}) = K^{(t+1)} w_{\mathrm{KQ}}(x_{1:N}^{(t+1)})$. 
And the solution is
\begin{align*}
    w^{(t+1)} = \frac{1}{1+\eta} w^{(t)} + \frac{\eta}{1+\eta} w_{\mathrm{KQ}}(x_{1:N}^{(t+1)}).
\end{align*}


\begin{theorem}[Monotone improvement over the initial kernel quadrature rule]
\label{thm:wift-kq-descent}
Let $\mu_N^{(t)}=\sum_{i=1}^N w_i^{(t)}\delta_{x_i^{(t)}}$ and $X_t=(x_1^{(t)},\ldots,x_N^{(t)})$. Define the MMD energy
$\mathcal{E}(X,w)=\frac12\operatorname{MMD}^2\left(\sum_{i=1}^N w_i\delta_{x_i},\pi\right)$,
and let $w_{\mathrm{KQ}}(X)=K(X)^{-1}m_\pi(X)$ denote the unconstrained kernel quadrature weights associated with the locations $X$.

Assume that $k$ is twice continuously differentiable, that $K(X_t)$ is positive definite for every $t$, and that the weights remain uniformly positive, namely $w_i^{(t)}\geq \underline w>0$ for every $i,t$. Assume moreover that, for every admissible $w$, the map $X\mapsto\mathcal{E}(X,w)$ has an $L_X$-Lipschitz gradient on a set containing all iterates. If $\tau\alpha\leq 2\underline w/L_X$, then
\begin{align}
\operatorname{MMD}
\left(
\mu_N^{(t+1)},\pi
\right)
\leq
\operatorname{MMD}
\left(
\mu_N^{(t)},\pi
\right)
\label{eq:wift-mmd-monotonicity}
\end{align}
for every $t\geq0$. 
\end{theorem}

\begin{proof}
For fixed $X$ and $w$, let
$g_i(X,w)=\sum_{j=1}^N w_j\nabla_2 k(x_j,x_i)-\nabla m_\pi(x_i)$.
Differentiating the MMD energy with respect to $x_i$ and using symmetry of the kernel gives
$\nabla_{x_i}\mathcal{E}(X,w)=w_i g_i(X,w)$.
Hence, the position update can be written as
$x_i^{(t+1)}=x_i^{(t)}-\frac{\tau\alpha}{w_i^{(t)}}\nabla_{x_i}\mathcal{E}(X_t,w^{(t)})$.
Thus, the Wasserstein substep is a diagonally preconditioned gradient step for $X\mapsto\mathcal{E}(X,w^{(t)})$.

By the $L_X$-smoothness of $\mathcal{E}(\cdot,w^{(t)})$, the descent lemma gives
\begin{align}
\mathcal{E}(X_{t+1},w^{(t)})
&\leq
\mathcal{E}(X_t,w^{(t)})
+
\sum_{i=1}^N
\left\langle
\nabla_{x_i}\mathcal{E}(X_t,w^{(t)}),
x_i^{(t+1)}-x_i^{(t)}
\right\rangle
\nonumber\\
&\qquad
+
\frac{L_X}{2}
\sum_{i=1}^N
\left\|
x_i^{(t+1)}-x_i^{(t)}
\right\|^2
\nonumber\\
&=
\mathcal{E}(X_t,w^{(t)})
-
\tau\alpha
\sum_{i=1}^N
\left(
w_i^{(t)}
-
\frac{L_X\tau\alpha}{2}
\right)
\left\|
g_i(X_t,w^{(t)})
\right\|^2.
\label{eq:wasserstein-descent}
\end{align}
Since $w_i^{(t)}\geq\underline w$ and $\tau\alpha\leq2\underline w/L_X$, we have
$w_i^{(t)}-L_X\tau\alpha/2\geq0$. Therefore,
$\mathcal{E}(X_{t+1},w^{(t)})\leq\mathcal{E}(X_t,w^{(t)})$.

We next consider the weight update. For fixed $X$, the dependence of the MMD energy on $w$ is quadratic:
$\mathcal{E}(X,w)=\frac12 w^\top K(X)w-m_\pi(X)^\top w+C_\pi$,
where $C_\pi$ is independent of $w$. Since $K(X)$ is positive definite, the unique minimizer is
$w_{\mathrm{KQ}}(X)=K(X)^{-1}m_\pi(X)$.
Completing the square gives
\begin{align}
\mathcal{E}(X,w)
=
\mathcal{E}\left(X,w_{\mathrm{KQ}}(X)\right)
+
\frac12
\left\|
w-w_{\mathrm{KQ}}(X)
\right\|_{K(X)}^2,
\label{eq:mmd-kq-decomposition}
\end{align}
where $\|v\|_{K(X)}^2=v^\top K(X)v$.

By the weight update,
$w^{(t+1)}-w_{\mathrm{KQ}}(X_{t+1})
=
(1+\eta)^{-1}
\left(
w^{(t)}-w_{\mathrm{KQ}}(X_{t+1})
\right)$.
Consequently,
\begin{align}
&
\mathcal{E}(X_{t+1},w^{(t+1)})
-
\mathcal{E}\left(X_{t+1},w_{\mathrm{KQ}}(X_{t+1})\right)
\nonumber\\
&\qquad=
\frac{1}{(1+\eta)^2}
\left[
\mathcal{E}(X_{t+1},w^{(t)})
-
\mathcal{E}\left(X_{t+1},w_{\mathrm{KQ}}(X_{t+1})\right)
\right].
\label{eq:kq-weight-contraction}
\end{align}
Since $(1+\eta)^{-2}\leq1$ and $w_{\mathrm{KQ}}(X_{t+1})$ minimizes
$w\mapsto\mathcal{E}(X_{t+1},w)$, it follows that
$\mathcal{E}(X_{t+1},w^{(t+1)})
\leq
\mathcal{E}(X_{t+1},w^{(t)})$.

Combining the two substeps yields
\begin{align}
\mathcal{E}(X_{t+1},w^{(t+1)})
\leq
\mathcal{E}(X_{t+1},w^{(t)})
\leq
\mathcal{E}(X_t,w^{(t)}).
\label{eq:full-wift-descent}
\end{align}
Hence, $\mathcal{E}(X_t,w^{(t)})$ is nonincreasing. Since
$\mathcal{E}(X,w)=\frac12\operatorname{MMD}^2(\sum_i w_i\delta_{x_i},\pi)$,
we obtain
$\operatorname{MMD}(\mu_N^{(t+1)},\pi)
\leq
\operatorname{MMD}(\mu_N^{(t)},\pi)$.

Finally, if $w^{(0)}=w_{\mathrm{KQ}}(X_0)$, monotonicity implies
$\mathcal{E}(X_t,w^{(t)})
\leq
\mathcal{E}(X_0,w^{(0)})
=
\mathcal{E}(X_0,w_{\mathrm{KQ}}(X_0))$.
Equivalently,
\begin{align}
\operatorname{MMD}^2
\left(
\mu_N^{(t)},\pi
\right)
\leq
\operatorname{MMD}^2
\left(
\sum_{i=1}^N
w_{\mathrm{KQ},i}(X_0)
\delta_{x_i^{(0)}},
\pi
\right).
\end{align}
Taking square roots proves \eqref{eq:wift-improves-initial-kq}.
\end{proof}


